II) \emph{Reducción al caso en que $A$ es un anillo local que es un dominio
entero de dimensión 1, $B$ un anillo de valuación discreta, clausura
integral de $A$, y el morfismo
$\operatorname{Spec}(B)\to\operatorname{Spec}(A)$ es radicial.} ---
Sea $K$ el cuerpo de fracciones de $A$. Se sabe (\textbf 0, 23.2.7) que $B$
es la intersección de una familia $(V_\lambda)$ de anillos de
valuación discreta tal que, para todo $x\in K$, se tiene
$x\in V_\lambda$ salvo para un número finito de índices $\lambda$.
Se tiene, pues, una sucesión exacta de $A$-módulos
\[
0\longrightarrow B\longrightarrow K\longrightarrow
\bigoplus_\lambda(K/V_\lambda).
\]
Pongamos $K'=K\otimes_A A'$, $V'_\lambda=V_\lambda\otimes_A A'$;
por planitud, se deduce de la sucesión exacta anterior una nueva
sucesión exacta
\[
0\longrightarrow B'\longrightarrow K'\longrightarrow
\bigoplus(K'_\lambda/V'_\lambda)
\]
y, por consiguiente, $B'=\bigcap_\lambda V'_\lambda$. Además,
$\operatorname{Spec}(K')$ es la fibra del morfismo
$\operatorname{Spec}(A')\to\operatorname{Spec}(A)$ sobre el punto
genérico de $\operatorname{Spec}(A)$; por tanto, $K'$ es un anillo
noetheriano \emph{normal} y, en consecuencia, un producto directo de un
número finito de dominios enteros integralmente cerrados. El producto
directo $L'$ de los cuerpos de fracciones de estos últimos es el
anillo total de fracciones de $A'$ y, por tanto, también el de $B'$.
Se puede, pues, aplicar el lema (6.14.1.1) y, si se demuestra que cada
$V'_\lambda$ es un anillo \emph{normal}, lo mismo ocurrirá con $B'$.
